\documentclass[11pt]{article}
\usepackage{cmap}
\usepackage[margin=1.1in]{geometry}
\usepackage{amsmath,amssymb,amsthm}
\usepackage[T1]{fontenc}
\usepackage{lmodern}
\input{glyphtounicode}
\pdfgentounicode=1
\usepackage{xcolor}
\usepackage{hyperref}
\definecolor{linkblue}{rgb}{0.1,0.2,0.55}
\hypersetup{colorlinks=true,linkcolor=linkblue,citecolor=linkblue,urlcolor=linkblue}
\emergencystretch=3em
\newtheorem{theorem}{Theorem}
\newtheorem{lemma}[theorem]{Lemma}
\newtheorem{corollary}[theorem]{Corollary}
\theoremstyle{remark}
\newtheorem{remark}[theorem]{Remark}
\newtheorem{problem}[theorem]{Open problem}

\title{A Weighted Tur\'an Unit-Step Inequality for Modified Bessel Functions\\and Its Determinant Corollary}
\author{Lluis Eriksson}
\date{8 November 2026 --- revision v2}
\providecommand{\SourceFreeze}{unreleased}

\begin{document}
\maketitle

\begin{abstract}
For $x > 0$ and real $m \ge 1$ define, with $I_\mu$ the modified
Bessel function of the first kind,
\[
\varphi_m(x) \;=\;
\frac{(m-1)\,I_{m-1}(x)^2 + (m+1)\,I_{m+1}(x)^2}{m\,I_m(x)^2}.
\]
We prove that $\varphi_m(x) < \varphi_{m+1}(x)$ for every $x>0$ and
every real $m \ge 1$ --- a weighted Tur\'an-type monotonicity
statement we have not found in the literature, although every
ingredient of the proof is classical. The proof is fully elementary
and uses exactly the arsenal of the companion note on unit-step
order-monotonicity: the three-term recurrence and the Amos-type
bound $\rho_m < x/(a + \sqrt{a^2+x^2})$, $a = m + \tfrac12$, for
$\rho_m = I_{m+1}/I_m$. The proof uses an exact factorization:
with $u = \rho_m(x)$, $c = 1/x$, $S = u + 1/u$, $P = 1/u - u$,
\[
\tfrac12\bigl(\varphi_{m+1}-\varphi_m\bigr)
\;=\; (S - 3c)\,\bigl(P - (2m+1)c\bigr) \;+\; (2m+1)\,c^2 ,
\]
an algebraic identity once $\rho_{m-1}$ and $\rho_{m+1}$ are
eliminated by the recurrence. The second factor is positive by the
calibrated Amos bound --- it is precisely the unit-step inequality of
the companion note --- and the first because the same bound forces
$u < x/(2m+1) \le x/3$. As an application we obtain the strict
ordering $c_{mn} < 0$ ($m<n$) of the $2\times2$ determinants of the
coefficient arrays of a $\pi$-local surface expansion arising in an
audit-first programme on two-dimensional $SU(2)$ lattice gauge
theory; this is the structural half of a positivity problem whose
remaining half --- a signed-series positivity question --- we state
precisely and leave open. No numerical cancellation rate is asserted
in this revision. An accompanying earlier Lean formalization encodes the ordered-field
implication with the analytic recurrences and Amos bound as named
hypotheses. Its archived verification record is supplied separately;
no new Lean or numerical run is claimed by this revision.
\end{abstract}

\section{Statement}

Throughout, $x > 0$, $I_\mu$ is the modified Bessel function of the
first kind, all values $I_\mu(x)$ for $\mu \ge 0$ are positive, and
$\rho_\mu(x) = I_{\mu+1}(x)/I_\mu(x) \in (0,1)$.

\begin{theorem}[weighted Tur\'an-type unit-step monotonicity]
\label{thm:main}
For every $x > 0$ and every real $m \ge 1$,
\begin{equation}
\label{eq:main}
\varphi_m(x) \;=\;
\frac{(m-1)\,I_{m-1}(x)^2 + (m+1)\,I_{m+1}(x)^2}{m\,I_m(x)^2}
\;<\; \varphi_{m+1}(x),
\end{equation}
strictly. Equivalently, in terms of the consecutive ratios
$\rho_\mu = I_{\mu+1}/I_\mu$:
\[
\frac1m\Bigl(\frac{m-1}{\rho_{m-1}^2} + (m+1)\rho_m^2\Bigr)
\;<\;
\frac1{m+1}\Bigl(\frac{m}{\rho_m^2} + (m+2)\rho_{m+1}^2\Bigr).
\]
\end{theorem}

The weights $(m-1, m, m+1)$ are not decorative: they arise from a
specific surface expansion (Section~\ref{sec:app}), and --- as the
proof shows --- this choice admits the exact factorization below over the classical Amos
bound. No uniqueness classification of weights is asserted here; compare
\cite{Amos, Segura2011, RuizSegura}. We have not found
\eqref{eq:main} stated in the Tur\'an-type literature for modified
Bessel functions \cite{TN, Baricz1010, BariczTuranians}, though it
sits squarely in its tradition: for fixed $x$, log-concavity of
$\mu \mapsto I_\mu(x)$ (the classical Tur\'an inequality
\cite{TN, Baricz1010}) controls the \emph{first-order} behaviour of
the ratio sequence, while \eqref{eq:main} is a statement about a
\emph{weighted second-order} combination.

\section{The proof}

We use two classical facts, valid for all real orders
(\cite[10.29.1]{DLMF}): the three-term recurrence, in the two forms
\begin{equation}
\label{eq:rec}
\frac{1}{\rho_{m-1}} = \rho_m + \frac{2m}{x},
\qquad
\rho_{m+1} = \frac{1}{\rho_m} - \frac{2(m+1)}{x},
\end{equation}
and the classical Amos-type upper bound \cite{Amos,Segura2011}.
Theorem~2 of Ruiz-Antol\'in and Segura \cite{RuizSegura}, applied
with their order $\nu=m+1\ge2$, supplies exactly the bound used here:
\begin{equation}
\label{eq:amos}
\rho_m(x) \;<\; U_m(x) := \frac{x}{a + s},
\qquad a = m + \tfrac12,\quad s = \sqrt{a^2 + x^2}.
\end{equation}

\subsection*{Step 1: elimination}

Write $u = \rho_m(x) \in (0,1)$ and $c = 1/x$. Substituting
\eqref{eq:rec} into the $\rho$-form of \eqref{eq:main} and clearing
the positive denominators $m, m+1$, the claim
$\varphi_{m+1} - \varphi_m > 0$ becomes $\Delta_m > 0$ where
\[
\Delta_m \;=\;
\frac{m^2}{u^2} - (m^2-1)\Bigl(u + 2mc\Bigr)^2
+ m(m+2)\Bigl(\frac1u - 2(m+1)c\Bigr)^2 - (m+1)^2 u^2 .
\]

\subsection*{Step 2: the exact factorization}

\begin{lemma}
\label{lem:factor}
With $S = u + 1/u$ and $P = 1/u - u$,
\[
\tfrac12\bigl(\varphi_{m+1}-\varphi_m\bigr)
\;=\; \frac{\Delta_m}{2m(m+1)}
\;=\; (S - 3c)\,\bigl(P - (2m+1)c\bigr) \;+\; (2m+1)\,c^2 .
\]
\end{lemma}

\begin{proof}
Expanding $\Delta_m$ and collecting, the coefficients of $1/u^2$ and
$u^2$ are $\pm 2m(m+1)$, the linear-in-$c$ terms are
$-4m(m+1)c\,[(m-1)u + (m+2)/u]$, and the $c^2$ terms sum to
$8m(m+1)(2m+1)c^2$; dividing by $2m(m+1)$,
\[
\frac{\Delta_m}{2m(m+1)} = \Bigl(\frac{1}{u^2} - u^2\Bigr)
- \frac{2}{x}\Bigl((m-1)u + \frac{m+2}{u}\Bigr)
+ \frac{4(2m+1)}{x^2}.
\]
Now $(m-1)u + (m+2)/u = \tfrac12[(2m+1)S + 3P]$ and
$1/u^2 - u^2 = PS$, so the right side equals
$PS - c(2m+1)S - 3cP + 4(2m+1)c^2
= (S-3c)(P-(2m+1)c) + (2m+1)c^2$.
\end{proof}

\subsection*{Step 3: positivity of both terms}

The second factor is strictly positive: $P - (2m+1)c > 0$ is
\emph{exactly} the unit-step order-monotonicity inequality
$1/\rho_m - \rho_m > (2m+1)/x$ of the companion note
\cite{Companion}, which follows from \eqref{eq:amos} by the exact
calibration $1/U_m - U_m = (2a)/x = (2m+1)/x$ together with the
strict decrease of $t \mapsto 1/t - t$ on $(0,\infty)$.

The first factor is strictly positive: since $s > a$ in
\eqref{eq:amos}, $U_m < x/(2a) = x/(2m+1)$, so
$u < x/(2m+1) \le x/3$ for $m \ge 1$, hence
$S > 1/u > (2m+1)/x \ge 3/x = 3c$.

The third term $(2m+1)c^2$ is positive. This proves
Theorem~\ref{thm:main}. \qed

\begin{remark}[the calibration, for the third time]
\label{rem:third}
The bound $U_m$ of \eqref{eq:amos} is the positive root of
$1/t - t = (2m+1)/x$: it is the extremal function of the threshold
that appears in the factorization. In the companion note
\cite{Companion} this exact calibration was the engine of the
unit-step monotonicity of $(\log I_\nu)'$ and was shown to be
\emph{equivalent} to the Amos bound; here the same calibrated
threshold reappears, multiplied against a second factor whose
positivity is again a consequence of the same bound. Within the
audit programme from which these notes emerge, this is the third
lock opened by the same key, which suggests the calibration
identity, rather than any particular inequality, is the structural
invariant of this family of problems.
\end{remark}

\section{Application: ordering of the $\pi$-local surface expansion}
\label{sec:app}

The following objects arise in an audit-first research programme on
the two-dimensional $SU(2)$ lattice gauge theory with Wilson action
(Griffiths-II positivity on $2$-complexes; see the programme records
\cite{MasterMap} for context). For $\beta > 0$ let
$r_k = I_{k+1}(\beta)/I_1(\beta)$, $k \ge 0$, and define the
coefficient arrays ($k \ge 1$)
\[
A_k = r_{k-1}^2 \!\!\sum_{j \in \{k-2,\,k\},\; j \ge 0}\!\! (j+1)\, r_j^2,
\qquad
B_k = k\, r_{k-1}^4 .
\]
A direct computation shows $A_k/B_k = \varphi_k(\beta)$ for every
$k \ge 1$ --- including $k = 1$, where the boundary truncation of the
sum (the $j = -1$ term is absent) is annihilated by the weight
$(k-1)$ in $\varphi_k$, so no separate edge case is needed.

\begin{corollary}[determinant ordering]
\label{cor:order}
For all $\beta > 0$ and all integers $1 \le m < n$,
\[
c_{mn} \;=\; A_m B_n - A_n B_m \;=\;
B_m B_n\,\bigl(\varphi_m(\beta) - \varphi_n(\beta)\bigr) \;<\; 0 .
\]
\end{corollary}

\begin{proof}
$B_k > 0$, and $\varphi_m < \varphi_n$ by Theorem~\ref{thm:main}
iterated over unit steps from $m$ to $n$.
\end{proof}

\begin{lemma}[absolute convergence of the signed application series]
\label{lem:convergence}
For fixed $\beta>0$, the series defining $W(t)$ below is absolutely
convergent, uniformly for real $t$.
\end{lemma}
\begin{proof}
The positive power series for integer-order $I_j$ \cite[10.25.2]{DLMF}
gives, for $j\ge0$,
\[
 I_j(\beta)=\sum_{h\ge0}\frac{(\beta/2)^{j+2h}}{h!(j+h)!}
 \le \frac{(\beta/2)^j}{j!}\exp(\beta^2/4).
\]
Thus $r_k$ has a factorial tail for fixed $\beta$. From their definitions,
$A_k$ and $B_k$ have finite polynomially weighted sums of every fixed
order; the finitely many small indices, including $k=1$, cause no problem.
Moreover $|c_{mn}|\le A_mB_n+A_nB_m$, while the absolute value of the
trigonometric bracket defining $W$ is at most $2(m+n)$. The corresponding
double sum is bounded by products of the finite zeroth and first moments
of $A_k$ and $B_k$. This gives a summable bound independent of $t$.
Absolute convergence justifies the parity rewriting below; it supplies no
sign for the resulting sum.
\end{proof}

\begin{problem}[what remains open, stated honestly]
\label{prob:open}
The positivity lemma that Corollary~\ref{cor:order} serves is: for
$t = \pi - \varepsilon$, $\varepsilon \in (0, 0.4]$, and all
$\beta > 0$,
\[
W(t) \;=\; \sum_{1 \le m < n} c_{mn}\,
\bigl[(m-n)\sin((m+n)t) + (m+n)\sin((n-m)t)\bigr] \;<\; 0 .
\]
Writing $p = n-m$, $q = n+m$, the substitution $t = \pi -
\varepsilon$ gives the \emph{exact identity}, valid for every pair,
\[
W = \sum_{m<n} |c_{mn}|\,(-1)^{m+n}
\bigl[q\sin(p\varepsilon) - p\sin(q\varepsilon)\bigr],
\]
where Corollary~\ref{cor:order} supplied $c_{mn} = -|c_{mn}|$. For
the pairs with $q\varepsilon \le \pi$ the bracket is strictly
positive (decrease of $\sin(y)/y$ on $(0,\pi]$); for
$q\varepsilon > \pi$ it is not sign-definite, so the sum is a
signed-weight competition. The claim $W<0$ is an
\emph{alternating-parity dominance} problem, exact in the rewriting and
open in the dominance. Lemma~\ref{lem:convergence} makes the displayed
series well defined for each fixed $\beta$; it does not establish its sign.
Earlier exploratory totals and precision prescriptions are not presented
here as validated numerical data: their parameter, truncation and error
specifications were incomplete. Cancellation may make naive summation
unreliable, but it excludes neither a sufficiently sharp estimate-based
proof nor alternative numerical methods. No universal decimal-precision
threshold is claimed. A future numerical study must specify
$(\beta,\varepsilon)$, cutoffs, working precision, tail bounds and error
control. The problem remains open. The coefficient arrays are the actual
proved application; their role in a gauge-theory programme is context,
not a derivation of a physical observable or a positivity theorem.
\end{problem}

\section{Machine verification}

The accompanying historical record reports a Lean~4/Mathlib check of
this ordered-field argument in \nolinkurl{PhiMonotone.lean}. The recorded
command is \texttt{lake env lean}, with Lean~\texttt{4.30.0-rc2} and
Mathlib revision
\begin{center}\small
\nolinkurl{cd3b69baae9cd81a572a3720f2372655eca39038}.
\end{center}
The recorded declarations are:
\begin{itemize}
\item \nolinkurl{amos_calibration'}: the exact calibration;
\item \nolinkurl{amos_small}: the inequality $u<x/(2m+1)$ from
\eqref{eq:amos};
\item \nolinkurl{phi_unit_step}: Lemma~\ref{lem:factor} and positivity;
\item \nolinkurl{phi_step_of_recurrences}: Theorem~\ref{thm:main}, with
\eqref{eq:rec} and \eqref{eq:amos} as the named hypotheses
\texttt{hrec1}, \texttt{hrec2}, \texttt{hAmos}.
\end{itemize}
The archived log, dated 9~July~2026, reports the following standard-axiom
output; this revision has not replayed it:
\begin{verbatim}
'amos_calibration'' depends on axioms:
    [propext, Classical.choice, Quot.sound]
'amos_small' depends on axioms:
    [propext, Classical.choice, Quot.sound]
'phi_unit_step' depends on axioms:
    [propext, Classical.choice, Quot.sound]
'phi_step_of_recurrences' depends on axioms:
    [propext, Classical.choice, Quot.sound]
\end{verbatim}
These are the archived run's reported exit code $0$ and absence of
\texttt{sorry} or \texttt{admit}, not new run results. As in the
companion note, the encoded statement is the ordered-field
implication [recurrences $+$ Amos bound $\Rightarrow$ monotonicity];
the analytic inputs themselves are classical and cited.

\section{Reproducibility and computational scope}

The supplied files are \nolinkurl{PhiMonotone.lean} and the unchanged
historical log \nolinkurl{LEAN-VERIFICATION-LOG.txt}. Their programme
source commit is
\begin{center}\small
\nolinkurl{871bb2bee8410dfb5b225d7961203eefaca5c15e}.
\end{center}
The source paths at that commit are
\begin{quote}\small
\nolinkurl{papers/phi-lemma/PhiMonotone.lean}\\
\nolinkurl{papers/phi-lemma/LEAN-VERIFICATION-LOG.txt}.
\end{quote}
The historical log records the command
\begin{quote}\small\ttfamily
lake env lean PhiMonotone.lean
\end{quote}
with Lean~\texttt{4.30.0-rc2} and the
full Mathlib revision displayed above. These are historical provenance,
not a newly authenticated replay in that environment. The package
\texttt{SOURCE207.json} binds the exact manuscript and supplied files by
Git-LF hashes; the build record identifies the frozen source commit,
actual Linux tool versions, process exits and exact PDF hash.

This revision makes no fresh symbolic or numerical verification claim.
Earlier finite-grid experiments do not prove the real-order theorem,
validate the open sign claim, or establish a universal working precision.
The analytic argument and Lemma~\ref{lem:convergence} are the evidence
provided here. A formal rerun, if performed, must use a clean pinned
Mathlib environment and retain complete exit and axiom logs. The recorded
conditional theorem is not a formal construction of analytic Bessel
functions or a complete gauge-theory argument.

\section{Model assessments and review context}

Codex assisted manuscript preparation and review orchestration. An AI
assessment is not human peer review, formal verification, a priority
ruling or editorial acceptance. Exact served backend, reasoning effort
and reviewer independence are not inferred from the product label. The
current PDF is frozen before any new exact-artifact numerical review is
issued; such a report is a separate hash-bound record, not a number
inherited from an earlier version. Full issued reports, objections and
actual review context must accompany any publication. The earlier v1
artifact and its review remain preserved; the present scientific changes
address attribution, convergence, computational scope and reproducibility.
No rating of v1 is claimed for this revision.

\subsection*{Acknowledgements}
This note is part of an audit-first research programme whose records,
Lean sources and numerical suites are public at
\url{https://github.com/lluiseriksson/THE-ERIKSSON-PROGRAMME}. The
statement is presented as a lemma with an application and an open
problem, in that order; nothing here is claimed about the open
problem beyond its precise reduction, and nothing is claimed about
four-dimensional gauge theory.

\begin{thebibliography}{10}\small

\bibitem{Amos}
D.~E. Amos,
\emph{Computation of modified Bessel functions and their ratios},
Math. Comp. \textbf{28} (1974), 239--251.

\bibitem{Segura2011}
J.~Segura,
\emph{Bounds for ratios of modified Bessel functions and associated
Tur\'an-type inequalities},
J. Math. Anal. Appl. \textbf{374} (2011), 516--528.

\bibitem{RuizSegura}
D.~Ruiz-Antol\'in and J.~Segura,
\emph{A new type of sharp bounds for ratios of modified Bessel
functions},
J. Math. Anal. Appl. \textbf{443} (2016), 1232--1246;
arXiv:1606.02008.

\bibitem{TN}
V.~R. Thiruvenkatachar and T.~S. Nanjundiah,
\emph{Inequalities concerning Bessel functions and orthogonal
polynomials},
Proc. Indian Acad. Sci. Sect. A \textbf{33} (1951), 373--384.

\bibitem{Baricz1010}
\'A.~Baricz and S.~Ponnusamy,
\emph{On Tur\'an type inequalities for modified Bessel functions},
Proc. Amer. Math. Soc. \textbf{141} (2013), 523--532;
arXiv:1010.3346.

\bibitem{BariczTuranians}
\'A.~Baricz,
\emph{Bounds for Tur\'anians of modified Bessel functions},
Expo. Math. \textbf{33} (2015), 223--251; arXiv:1202.4853.

\bibitem{DLMF}
F.~W.~J. Olver et al.,
\emph{NIST Digital Library of Mathematical Functions},
\S\S10.25 and 10.29, \url{https://dlmf.nist.gov/10.25},
\url{https://dlmf.nist.gov/10.29}.

\bibitem{Companion}
L.~Eriksson,
\emph{An elementary recurrence--Amos proof of the unit-step
order-monotonicity of $(\log I_\nu)'$, with a Feynman--Hellmann
application to two-dimensional lattice gauge theory},
July 2026;
\url{https://github.com/lluiseriksson/THE-ERIKSSON-PROGRAMME/tree/main/papers/bessel-amos-fh}.

\bibitem{MasterMap}
L.~Eriksson,
\emph{THE MASTER MAP: an audit-first navigation guide to the
conditional construction of 4D $SU(N)$ Yang--Mills with mass gap},
ai.viXra:2602.0096, and the programme repository
\url{https://github.com/lluiseriksson/THE-ERIKSSON-PROGRAMME}.

\end{thebibliography}

\end{document}
